\section*{Declaración de Uso de Herramientas de Inteligencia Artificial}
\addcontentsline{toc}{section}{Declaración de Uso de Herramientas de Inteligencia Artificial}

% ============================================================================
% REVISAR ANTES DE FIRMAR. Esta es una declaracion formal y debe reflejar
% exactamente lo que hizo el grupo. Ajustar el alcance, los integrantes que
% usaron la herramienta y las fechas antes de la entrega.
% ============================================================================

\subsection*{1. Identificación de la actividad}
\begin{itemize}
    \item \textbf{Curso:} Arquitectura de Computadores
    \item \textbf{Código del curso:} 433
    \item \textbf{Nombre de la actividad:} Proyecto 2: Juego en Assembly sobre Pochoco SoC
    \item \textbf{Fecha de entrega:} Martes 29 de septiembre de 2026
    \item \textbf{Nombre(s) de(los) estudiante(s):} Josefa Alcaíno Lira, Felipe Izcúe Ferrer, Martín Norton Haigis
\end{itemize}

\subsection*{2. Declaración sobre uso de IA}
\noindent $\square$ No se utilizaron herramientas de inteligencia artificial en ninguna etapa del desarrollo de esta actividad.

\noindent $\boxtimes$ Sí se utilizaron herramientas de inteligencia artificial, según se detalla en la sección 3 de este documento.

\subsection*{3. Descripción de uso de IA}
\begin{table}[H]
    \centering
    \begin{tabularx}{\linewidth}{|l|X|}
        \hline
        \textbf{Herramienta empleada} & Claude (Anthropic), a través de la interfaz Claude Code. \\
        \hline
        \textbf{Fechas de uso} & 27 y 28 de septiembre de 2026. \\
        \hline
        \textbf{Estudiantes} & Martín Norton Haigis. \\
        \hline
        \textbf{Propósito del uso} &
        Desarrollo, con la herramienta, del código del proyecto a partir del enunciado
        y del repositorio base: el contador de ciclos y el módulo \code{game\_top} en
        Verilog, el programa del juego en assembly, el assembler en Python, sus pruebas
        automáticas, los testbenches de simulación y el Makefile.
        Diseño y ejecución de la verificación: simulación en Icarus Verilog, medición
        de ciclos por instrucción, modelo del generador pseudoaleatorio y síntesis con
        Yosys y nextpnr.
        Redacción de este informe y de la documentación del repositorio. \\
        \hline
        \textbf{Descripción de integración} &
        Las decisiones de diseño se tomaron durante el desarrollo y quedaron
        registradas en el código y en este informe. Todo el código se verificó con las
        pruebas y simulaciones descritas en la sección~\ref{sec:verificacion}, cuyos
        resultados se reportan tal como se obtuvieron. El grupo revisó el código y el
        texto, y comprende y puede explicar cada parte del diseño, del programa y del
        assembler. \\
        \hline
    \end{tabularx}
\end{table}

\subsection*{4. Declaración de responsabilidad académica}
\noindent Los firmantes de este documento declaramos que:
\begin{itemize}
    \item Hemos revisado críticamente todo el contenido entregado.
    \item Comprendemos plenamente el trabajo presentado y somos capaces de explicarlo y defenderlo académicamente.
    \item Asumimos responsabilidad por la calidad, veracidad y originalidad del contenido entregado, independientemente del uso o no uso de herramientas de inteligencia artificial.
    \item La información declarada en este documento es veraz.
\end{itemize}

\vspace{1.5cm}
\noindent
\begin{tabular}{p{4.6cm}p{4.6cm}p{4.6cm}}
    \vspace{2cm} & \vspace{2cm} & \vspace{2cm} \\
    \hrulefill & \hrulefill & \hrulefill \\
    Josefa Alcaíno Lira & Felipe Izcúe Ferrer & Martín Norton Haigis \\
\end{tabular}
